\documentclass[12pt,a4paper]{letter}
\usepackage[a4paper,margin=1.7cm]{geometry}
\usepackage{siunitx}
\usepackage[colorlinks=true,citecolor=blue,linkcolor=blue,urlcolor=blue]{hyperref}
\sisetup{mode = match}
\setlength{\parskip}{0.35em}
\setlength{\medskipamount}{0.35em}
\address{%
P. Djorw\'e\\
Department of Physics, University of Ngaoundere, Cameroon;\\
STIAS, Stellenbosch University, South Africa\\
\texttt{djorwepp@gmail.com}}
\signature{P. Djorw\'e\\on behalf of all authors}
\date{\today}

\begin{document}
\begin{letter}{The Editors\\\textit{Physical Review A}}
\opening{Dear Editors,}

We submit the manuscript ``Finite-size and loss-profile limits on topological protection of Wigner negativity in a chiral Floquet SSH chain'' for consideration as a Regular Article in \textit{Physical Review A}.

Topological classifications are often invoked to argue that quantum resources inherit robustness, yet whether a bulk invariant constrains a non-Gaussian resource under dissipation, nonlinearity, and finite size is rarely tested directly. We ask where the boundary of topological protection lies for the local Wigner negativity of a two-step chiral Floquet SSH chain with Kerr nonlinearity, edge-selective loss, and a non-Gaussian cat state. The test compares the matched negativity contrast between two sectors with distinct Floquet invariants under identical conditions, and asks whether the contrast tracks the bulk invariant, the edge-mode localization, or neither.

The answer is neither. The contrast changes sign with system size under the stated \num{10}-trajectory estimator: $+\num{0.017}$ at two unit cells and $-\num{0.0018}$ to $-\num{0.0026}$ at three to five cells, with at least \num{4000} matched disorder realizations at each of the first three sizes. The exact master equation resolves the sign pattern at two, three, and four cells ($-\num{0.00032}\pm\num{0.00011}$ at four cells, \num{200} trajectories) and leaves five cells unresolved; finite-trajectory bias dominates the estimator's negative branch. A five-cell value predicted before measurement lies inside its pre-specified interval in all \num{4} conditions. A Kerr-free stratum preserves the sign pattern, so the nonlinearity amplifies the contrast but does not create it. Displaced Fock-state witnesses detect the two-cell signal without a phase-space grid (\qty{82.4}{\percent} versus \qty{10.8}{\percent} of disorder draws). The contrast vanishes across the topological transition, which shows that it does not track the bulk invariant. Detection margins over the minimum detectable effect are modest ($\approx\num{2.6}\times$ at eight modes, $\approx\num{1.3}\times$ at six modes), and magnitudes are not converged with respect to the Fock cutoff.

The manuscript separates three claims that are often conflated: a bulk invariant, a finite-chain edge mode, and the persistence of a resource. Here the resource is detectable and its finite-size trend is predictable, but it is not protected. We claim no topological protection, no metrological gain, and no experimental validation. The result gives a concrete boundary condition for future studies of nonclassical resources in open Floquet chains, which suits the readership of \textit{Physical Review A}.

All parameters, paired data, and reproducibility materials accompany the submission. The work is original and is not under consideration elsewhere. There is no external funding and there are no competing interests.

\closing{Sincerely,}
\end{letter}
\end{document}
